\documentclass[11pt]{article}

\usepackage{amsmath}
\usepackage{amssymb}
\usepackage{amsthm}
\usepackage{array}
\usepackage{booktabs}
\usepackage[margin=1in]{geometry}

\newtheorem{theorem}{Theorem}
\newtheorem{lemma}[theorem]{Lemma}
\newtheorem{corollary}[theorem]{Corollary}
\newtheorem{proposition}[theorem]{Proposition}
\theoremstyle{definition}
\newtheorem{definition}[theorem]{Definition}
\theoremstyle{remark}
\newtheorem{remark}[theorem]{Remark}

\newcommand{\dist}{\operatorname{dist}}
\newcommand{\Ball}{B}

\title{Disproof of Conjecture 00000001021:\\
The Minimum Order of a $(3,6)$-Regular Girth-$12$ Tanner Graph Is Not $132$}
\author{earthking11}
\date{2026-09-13}

\begin{document}
\maketitle

\begin{abstract}
The conjecture asserts that the minimum number of vertices of a
$(3,6)$-regular Tanner graph of girth $12$ is $132$, and is described as
``computationally verifiable''. We show this is false. A Tanner graph is
bipartite with a degree-$3$ vertex class (variable nodes) and a degree-$6$
vertex class (check nodes). If the girth is at least $12$, then two distinct
paths of lengths $i,j\le 5$ from a fixed variable node cannot end at the same
vertex, since they would close a cycle of length $i+j\le 10<12$; hence the
radius-$5$ ball around any variable node is a tree. Its BFS layers are forced
to be
\[
1,\ 3,\ 15,\ 30,\ 150,\ 300,
\]
whose sum is $499$. Therefore every such graph has at least $499$ vertices,
and in particular none has $132$. Imposing the balance relation $3V=6C$ of a
$(3,6)$-regular bipartite graph strengthens the bound to $999$. Inspecting the
whole $(d_v,d_c)\in[2,8]^2$, even-girth-$4$--$16$ parameter family, no
Moore-type bound equals $132$; the girth-$12$ lower bound is either $499$
(tree-ball) or $999$ (balanced), both far above $132$.
\end{abstract}

\section{Definitions}

\begin{definition}[Tanner graph]
A \emph{Tanner graph} is a finite bipartite graph $G=(X\cup Y,E)$ in which
every edge joins a \emph{variable node} in $X$ to a \emph{check node} in $Y$.
It is \emph{$(d_v,d_c)$-regular} if every vertex of $X$ has degree $d_v$ and
every vertex of $Y$ has degree $d_c$. Thus a $(3,6)$-regular Tanner graph has
variable nodes of degree $3$ and check nodes of degree $6$.
\end{definition}

\begin{definition}[Girth]
The \emph{girth} of $G$ is the length of a shortest cycle in $G$. Since $G$ is
bipartite, all cycles have even length.
\end{definition}

We write $V=|X|$, $C=|Y|$, and $N=V+C$ for the total number of vertices.
Because $G$ is $(3,6)$-regular, double counting the edges gives
\begin{equation}\label{eq:balance}
    3V=6C, \qquad\text{equivalently}\qquad V=2C,\qquad N=3C=\tfrac32 V.
\end{equation}

\section{The collision lemma}

Fix a variable node $v\in X$ and let $\dist(v,\cdot)$ be graph distance in $G$.
Let $\Ball(v,r)=\{w : \dist(v,w)\le r\}$.

\begin{lemma}[No collisions up to radius $5$]\label{lem:collision}
Let $G$ be a Tanner graph of girth $g\ge 12$ and let $v$ be any vertex. Write
$r=\lfloor (g-1)/2\rfloor$. Then the ball $\Ball(v,r)$ is a tree. In
particular, for $g=12$, distinct paths from $v$ of lengths $i,j\le 5$ have
distinct endpoints.
\end{lemma}

\begin{proof}
Suppose two distinct paths $P_1,P_2$ from $v$ of lengths $i$ and $j$ share
their terminal vertex $w$. Walking along $P_1$ from $v$ to $w$ and back along
$P_2$ from $w$ to $v$ yields a closed walk of length $i+j$. If $P_1,P_2$ are
distinct, this closed walk contains a cycle; its length is at most $i+j$ and,
being a cycle in a bipartite graph, it is even and hence at least $2$. More
precisely, the symmetric difference of the two paths contains a cycle of
length $\le i+j$. If $i,j\le 5$ then $i+j\le 10<12\le g$, contradicting the
definition of girth. Hence no two such paths collide.

It remains to see that the whole ball, not just its endpoints, is a tree.
\begin{enumerate}
\item \emph{No edges inside a layer.} In a bipartite graph an edge joining two
vertices at the same distance from $v$ would create an odd cycle, impossible.
\item \emph{No two layers meet.} Distance is single-valued: a vertex cannot
have two different distances from $v$.
\item \emph{No cross edges.} An edge between a vertex at depth $d$ and a
vertex at depth $d'$ with $|d-d'|\neq 1$ would shortcut the BFS, forcing
$\dist(v,\cdot)\le\min(d,d')+1<\max(d,d')$, contradicting single-valuedness of
the other depth.
\end{enumerate}
Combining these facts, the edges of $\Ball(v,r)$ are exactly the tree edges
of BFS, with no extra edge possible among depths $0,\dots,r$ when
$r\le g/2-1$. For $g\ge 12$ we have $\lfloor (g-1)/2\rfloor\ge 5$, so the
radius-$5$ ball is a tree.
\end{proof}

\begin{remark}
The same argument shows that $\Ball(v,5)$ is a subgraph of $G$ with $499$
distinct vertices; it need not be an induced subgraph, but the collision
argument above already rules out every extra edge that could occur inside the
ball at these depths.
\end{remark}

\section{Forced BFS layers and the $499$ bound}

Let $G$ be $(3,6)$-regular of girth $g\ge 12$ and let $v$ be a variable node.
By Lemma~\ref{lem:collision}, the radius-$5$ ball around $v$ is a tree, so the
BFS layers are forced by the degrees, with one edge from each non-root vertex
used by its parent.

\begin{center}
\begin{tabular}{c l c c c}
\toprule
depth $d$ & vertex class at depth $d$ & children per vertex & layer size & cumulative\\
\midrule
$0$ & variable (root, no parent)          & $3$   & $1$   & $1$\\
$1$ & check ($6$ neighbours, $1$ parent)  & $5$   & $3$   & $4$\\
$2$ & variable ($3$ neighbours, $1$ parent)& $2$   & $15$  & $19$\\
$3$ & check                               & $5$   & $30$  & $49$\\
$4$ & variable                            & $2$   & $150$ & $199$\\
$5$ & check                               & ---   & $300$ & $499$\\
\bottomrule
\end{tabular}
\end{center}

The layer sizes are obtained by the recurrence
\[
\begin{aligned}
\ell_0&=1, &\quad \ell_1&=3\ell_0=3, &\quad \ell_2&=5\ell_1=15,\\
\ell_3&=2\ell_2=30, &\quad \ell_4&=5\ell_3=150, &\quad \ell_5&=2\ell_4=300.
\end{aligned}
\]

\begin{theorem}\label{thm:499}
Every $(3,6)$-regular Tanner graph of girth $12$ has at least $499$ vertices:
\[
|\Ball(v,5)|=1+3+15+30+150+300=499\le N .
\]
Consequently the minimum number of vertices of such a graph is not $132$, and
indeed no such graph has $132$ vertices.
\end{theorem}
\begin{proof}
Lemma~\ref{lem:collision} makes the six layers pairwise disjoint, and the
layer sizes are forced as above. Their sum is $499$, all inside $G$. Since
$132<499$, the claimed value cannot occur.
\end{proof}

Note that $499$ is only a lower bound on the true minimum; the disproof needs
nothing more.

\section{The balanced variant: $999$}

The count above ignores the balance relation~\eqref{eq:balance}, which a
Tanner graph must satisfy. Splitting the forced ball by vertex class: the
variable nodes at even depth number $1+15+150=166$, and the check nodes at odd
depth number $3+30+300=333$. Hence
\[
V\ge 166,\qquad C\ge 333 .
\]
By~\eqref{eq:balance}, $V=2C$, so $C\ge 333$ already gives $V\ge 666$ and
\[
N=V+C=3C\ge 3\cdot 333=999 .
\]
Equivalently, among the $333$ forced check nodes each is adjacent only to
variable nodes, and a $(3,6)$-regular graph needs twice as many variable as
check nodes; the check-side count is the binding constraint.

\begin{theorem}
Every $(3,6)$-regular Tanner graph of girth $12$ has at least $999$ vertices.
\end{theorem}

Both bounds ($499$ from the tree-ball alone, $999$ after balancing) are far
from $132$.

\section{No convention in the parameter family yields $132$}

For general $(d_v,d_c)$-regular Tanner graphs of even girth $2r$, the same
argument forces the BFS layers
\[
\ell_0=1,\qquad \ell_{i+1}=\ell_i\, b_i,
\]
where the branching is
\[
b_0=d_v,\qquad b_i=d_c-1\ \ (i\ \text{odd}),\qquad
b_i=d_v-1\ \ (i>0\ \text{even}).
\]
The tree-ball lower bound is $\sum_{i=0}^{r}\ell_i$. Combining it with the
balance relation $d_vV=d_cC$ gives a balanced lower bound on
$N=V+C$, obtained by requiring
$V\ge\sum_{i\ \text{even}}\ell_i$ and
$C=\tfrac{d_v}{d_c}V\ge\sum_{i\ \text{odd}}\ell_i$.

Scanning all $2\le d_v,d_c\le 8$ and all even girths
$4,6,\dots,16$ (the verification is in \texttt{reproduce.py}) yields
\emph{no} parameter choice whose tree-ball or balanced bound equals $132$.

For the specific family $(d_v,d_c)=(3,6)$ the rigorous balanced bounds are:
\begin{center}
\begin{tabular}{c c c}
\toprule
girth & tree-ball lower bound & balanced lower bound \\
\midrule
$4$  & $4$   & $9$\\
$6$  & $19$  & $24$\\
$8$  & $49$  & $99$\\
$10$ & $199$ & $249$\\
$12$ & $499$ & $999$\\
$14$ & $1999$ & $2499$\\
$16$ & $4999$ & $9999$\\
\bottomrule
\end{tabular}
\end{center}
The value $132$ matches none of these: the rigorous bounds are $24$ for girth
$6$, $99$ for girth $8$, $249$ for girth $10$, and $999$ for girth $12$.
Even the raw tree-ball bound for girth $12$ is $499$. There is therefore no
reading of ``girth $12$'' under which $132$ is the minimum: $132$ is simply
too small, and the claimed ``minimal graph counting'' is incorrect.

\section{Conclusion}

Conjecture 00000001021 is false. The radius-$5$ ball around a degree-$3$
variable node in a $(3,6)$-regular girth-$12$ Tanner graph is a tree with
$1+3+15+30+150+300=499$ vertices, so every such graph has at least $499$
vertices; the balanced bound is $999$. No $(d_v,d_c)$-regular girth-$12$
graph, and no member of the scanned parameter family, can have only $132$
vertices.

\end{document}
